% NAF 5166, batch 3 — Gallica views 41–60.
% Pass: Opus 5.5 (claude-opus-5-5), 2026-09-23 — first pass, unchecked against the views by a human.
%
% What the views are. Greyscale scans, portrait, about 4700 × 6450, one
% opening per view. Old leaves and half-leaves of several sizes, torn at the
% edges, hinged on modern guard sheets, some mounted high and some low on
% their guard; one guard (v. 55) carries two slips. Every view was first
% looked at as a 1000-px thumbnail, and every apparently blank view again at
% stretched contrast; the written views were read as full-width bands at
% 2400 px cut over the leaf, with tight crops (1200–2400 px) for struck
% words, participles, numbers and every doubtful formula. Read only from the
% local mirror; nothing was fetched from Gallica. \page{N} is always the view.
%
% What the twenty views hold.
%
%   v. 41     Leaf « 12. ». The END of a piece begun before this batch: the
%             page opens mid-sentence (« à l'equation que nous venons de
%             trouver, car l'integrale de ρ dμ etant… »), so v. 41 CONTINUES
%             a text begun before it. Article 10: from V = 4πa³/3r + u
%             (n.o 2), V/2 + a dV/dr = −2πa²/3 at the surface, « Comme M.
%             Laplace l'a trouvé » (margin), a struck sentence calling it
%             « l'equation de M. Laplace … demontrée rigoureusement », and the
%             remark that the term 2πa²y is destroyed by the integral of
%             a³(ρ/2 + a dρ/dr) y' dϖ' dθ' sin θ' (n.o 4), whose factor is
%             identically nil. Ends « identiquement nul. » over a long rule;
%             the lower quarter is blank. Same subject, notation (ϖ, the V(
%             radical) and hand as batch 1's « Eclaircissement d'une
%             difficulté singuliere… » (v. 17–19); whether it is that memoir's
%             last leaf cannot be told from this batch alone.
%   v. 42     Verso of 12, blank (show-through). Not transcribed.
%   v. 43     Leaf « 13 », a smaller leaf: a single heading « Notes de M.r
%             Dela Grange » in a fine, careful hand unlike the text of v. 41;
%             rest blank. Apparently a wrapper title for what follows; what
%             it covers is not written on it.
%   v. 44     Verso of 13, blank. Not transcribed.
%   v. 45     Leaf « 14 ». PIECE: a short note without title, in the hand of
%             v. 41 — the plane Ax + By + Cz = 0 and the sphere about
%             (a, 0, 0): foot of the perpendicular, r = Aa/√(A²+B²+C²),
%             « r/a est le sinus de l'angle que le plan fait avec l'axe des
%             x », and its cosine reading. Upper two-thirds only.
%   v. 46     Verso of 14, blank. Not transcribed.
%   v. 47     Slip « 15 » « A ». PIECE (hand of v. 41), headed « page 623 »:
%             « Je mets x à la place de b je differentie », reduction to
%             −½ (1+x)^{−3/2} dx / x, expansion, term-by-term integration,
%             « + K ». Runs on to v. 49.
%   v. 48     Verso of 15, blank. Not transcribed.
%   v. 49     Slip « 16 » « B ». Continues v. 47: the constant K = −1 + l2
%             found by letting b → 0, and « la serie cherchée ». Ends at the
%             foot of the slip on « où l2 − ½lb = l(2/√b) ».
%   v. 50     Verso of 16, blank. Not transcribed.
%   v. 51     Leaf « 17 ». PIECE (hand of v. 41): arc of the ellipse,
%             x = a sin φ, y = b cos φ, reduced to E(c, φ) and « en general »
%             ∫√(dx²+dy²) = aE(c/a, φ). Ends two-thirds down.
%   v. 52     Verso of 17, blank. Not transcribed.
%   v. 53     Leaf « 18 ». SECOND HAND (careful, rounded, fine pen, d with a
%             straight stem, z like a 2): the calculation of v. 45 again, in
%             fuller prose — dz, the differential of the distance, the plane
%             through the origin perpendicular to the line with direction
%             A, B, C, « r/a est le sinus de l'angle… ». Ends two-thirds down.
%   v. 54     Verso of 18, blank. Not transcribed.
%   v. 55     One guard, two slips.
%             Upper slip « 19 », SECOND HAND, headed « Pag. 623 » (the page
%             cited at v. 47): expansions of log((1+√(1+b))/√b),
%             log(1+√(1+b)), −(1+b)^{−1/2}, log 2 = 1 − ½ + ⅓ − … — the
%             series of v. 47–49 by another route. Lower half of slip blank.
%             Lower slip « 20 »: a scrap of numerical reckoning in ink, no
%             words — sexagesimal subtractions, a doubled crescent like the
%             moon sign, and a small table « SC. », « ST. », « S.DAB » with
%             numbers shaped like tabular log-sines. Only the clear blocks are
%             given; the rest is described in a note.
%   v. 56     Back of the two slips. The upper is blank; the lower (back of
%             « 20 ») carries pencil tables of figures, part upside down, and
%             an ink group turned sideways. Given as a note only.
%   v. 57     Modern guard sheet, blank. Not transcribed.
%   v. 58     Leaf « 21 ». NEW PIECE, begins at the head of the leaf: a
%             first-person text to a « vous », without salutation, date or
%             signature — a letter or its draft. The writer thanks the
%             correspondent for obtaining a judgement on their work, is
%             surprised by « l'equation que vous m'indiquez d'après M.r De la
%             Grange », quotes two fourth-order equations in z, t, x, y, admits
%             « mon erreur … reconnue par des juges », and cites « la note »
%             and « la mécanique analytique » (1re édition). Runs on to v. 59.
%   v. 59     The sheet photographed open: left page = verso of leaf 21, which
%             continues v. 58 without a break (« ne convient-il donc pas de la
%             differencier quatre fois… »): the equation of the stretched
%             surface X(d²z/dx² + d²z/dy²) = 2Z, « la même que celle donnée
%             par M.r La Grange lui même p. 151 de la 1ère édition », and
%             the objection, reported from the correspondent, that the
%             equation « M.r De la Grande » found on the writer's hypothesis
%             contains that of the elastic lamina. The right-hand side of the
%             view is the left part of leaf 22's recto, transcribed at v. 60.
%   v. 60     Leaf « 22 », recto: continues v. 59 (heavily revised top
%             third); 1/r + 1/r′ in place of 1/r in the equation of the
%             lamina; « je sais bien que ce n'est pas ainsi qu'a du operer
%             M.r De la Grange »; ends « …ma parfaite soumission aux oracles
%             … jugement d'un autre grand géomètre. » with a full stop, the
%             rest of the page blank.
%
%   Where the batch begins and ends. View 41 continues a text begun before
%   it (it opens mid-sentence, and its article 10 cites n.os 2 and 4). View
%   60 closes its sentence and its text with a full stop and blank paper
%   below, with no signature; the verso of leaf 22 (presumably v. 61) was not
%   seen, so whether anything follows on it cannot be told here. The piece
%   of v. 58–60 is continuous across its three views.
%
%   The hands. (1) v. 41, 45, 47, 49, 51: the quick sloping cursive of batch
%   1 — looped d like a ∂, « est » with a long s, heavy ink, radical as a
%   large V( ; attributed to nobody here (the wrapper, v. 15, calls the
%   papers « autographes » of Lagrange). (2) v. 53 and the upper slip of v.
%   55: a careful rounded hand, fine pen, letters often detached, d with a
%   straight stem, z shaped like a 2, 5 shaped like ſ, 2 shaped like z. (3)
%   v. 58–60: a regular sloping hand with marked thick and thin strokes and
%   long horizontal exit strokes; it shares the straight d and the ſ-shaped
%   5 of (2) and may be the same hand written faster — not settled. The
%   writer of (3) writes of themself « égarée », « persuadée », « revenue »,
%   « frappée » (v. 58, 60), with the -ée plain; « étonné », « entraîné »
%   (v. 58) end in an ambiguous stroke. The wrapper (v. 15) says the volume
%   holds « notes etc de Mlle Sophie Germain »; this pass attributes none of
%   these leaves to anyone. (4) v. 43's heading: a fine careful hand,
%   ornate capitals — possibly (2), not settled. (5) the scrap « 20 »
%   (v. 55–56): figures only, the 4 written like an A.
%
%   Notation, fixed for the batch. The hand-(1) d, looped like a ∂, is set d,
%   as in batch 1. The angle written as a small π with curled bar is set ϖ;
%   the constant of 4πa³/3r, a larger Π-shaped letter, is set π. The looped
%   letter of « ρ dμ » is set ρ (v. 41). The angle of v. 51 is set φ. The
%   radical V( is set \sqrt{( )}. Logarithms are « l » (hand 1) and « log. »
%   (hand 2), as written. Products of factorial type keep their points
%   (3.5.7 / 2.4.6). Exponents written on the line (« b4 ») are set as
%   exponents and said in a note. No long s.
%
%   Numbers on the leaves, and which kind.
%   - « 12. » (v. 41), « 13 » (v. 43), « 14 » (v. 45), « 15 » (v. 47),
%     « 16 » (v. 49), « 17 » (v. 51), « 18 » (v. 53), « 19 » and « 20 »
%     (the two slips of v. 55), « 21 » (v. 58), « 22 » (v. 60): ink, top
%     right of each leaf or slip, one per physical piece, on the written
%     side only, never on a verso, running 12 → 22 without a gap across
%     four different hands of text. They continue the series « A. », « 1 »,
%     « 2 » of batch 1 and answer the library's certificate (v. 13: 26
%     leaves plus A and 6 bis), so they behave as the volume's foliation;
%     but they are in ink, in a figure unlike any of the text hands, and,
%     following the NAF 5161 practice and batch 1, no \folio{} is written.
%     (A human eye on the leaves would settle ink vs. pencil.)
%   - « A » (v. 47) and « B » (v. 49): capitals beside « 15 » and « 16 »,
%     in a heavier pen, pairing the two slips of one calculation; a
%     sequence mark, whose hand is not settled.
%   - « page 623 » (v. 47) and « Pag. 623 » (v. 55): a page reference in
%     the text, to a work not named on the leaves.
%   - In the text: article « 10. », « (n.o 2) », « (n.o 4) » (v. 41) — the
%     writer's article numbers; « p. 53 1re édition » and « p. 151 de la
%     1ère édition » (v. 58, 59; both doubtful) — page references to « la
%     mécanique analytique » and to La Grange.
%   - On the scrap « 20 » (v. 55–56): figures only, degrees and minutes and
%     log-like numbers; none is a numbering.
%   - v. 42, 44: a bold oblique ink mark like « 11 » or a paraph on the
%     corner of a sheet showing under the leaf; not on any leaf of this
%     batch; meaning not settled.
%
%   Apparatus. \ill{}: v. 47 ×1 (struck « const. = … »), v. 55 ×2 (a blotted
%   figure on the scrap; one digit of a log), v. 56 ×1 (one pencil digit),
%   v. 59 ×1 (a struck word), v. 60 ×3 (a word under a rature; two struck
%   words). No \folio{}.

% Coordinator's reconciliation (2026-09-23): the verso of leaf 22 is
%   view 61 (batch 4). Batch 4 read, mirrored through that leaf, phrases of
%   v. 58–60 (« M^r De la Grange », « frappée à la vue de son équation »,
%   a last line ending « géomètre. »), so v. 60 and v. 61 are the two sides
%   of one leaf; v. 61 carries a separate, unfinished draft. Which side
%   was written first is not settled. The numbers 12–22 are the library's
%   series (see batch 1's reconciliation), and no \folio{} is written.

\documentclass[11pt,a4paper]{article}
% The preamble shared by every transcription of the mathematicians' manuscripts in Gallica.
%
% It rests on one idea: **uncertainty compounds, it does not resolve.** A
% transcription that smooths over a doubtful word has destroyed the only thing
% it was made for — a reader must be able to tell, at every point, what was on
% the page and what was supplied. The apparatus macros below are therefore the
% critical apparatus, not decoration.
%
% The same commands are recognised by `scripts/render.mjs`, which produces the
% site's reading view, and by `scripts/tei.mjs`. A macro added here must be
% added there too, or rendering fails loudly — which is the wanted behaviour.
%
% Comments here are English, like the rest of the repository. The documents
% this preamble sets are French, and that is deliberate: see the skills.

\ProvidesPackage{germain}[2026/09/15 Transcriptions of mathematicians' manuscripts in Gallica]

\usepackage{amsmath,amssymb,amsthm}
\usepackage{mathrsfs}
% Kept from the parent preamble so the renderer's diagram support stays
% available; these manuscripts draw no commutative diagrams, and nothing here uses it.
\usepackage{tikz-cd}
\usepackage[english,french]{babel}
\usepackage{fontspec}

% --- Greek ------------------------------------------------------------------
% The text face stays fontspec's default, Latin Modern, which has no Greek:
% XeTeX drops every glyph it lacks with a line in the log and nothing on the
% page, and the Greek words the manuscripts quote vanished from the PDFs. So
% the Greek and Greek Extended blocks get a XeTeX character class of their own,
% and entering or leaving a run of it switches the family to CMU Serif — the
% same Computer Modern design, with polytonic Greek — and back. Only the
% family changes: a Greek word in \textit or \textbf keeps its shape and series.
% The transitions to and from every other class carry the tokens that class
% has towards an ordinary letter, so French spacing before « ; » or inside
% « » still holds after a Greek word. No grouping, so no brace or \endgroup
% can come out unbalanced; maths is left alone, since XeTeX inserts nothing
% there. Kept identical in `exercices.sty`.
\makeatletter
\newfontfamily\germainGreekfont{cmun}[
  Extension=.otf, NFSSFamily=germaingreek,
  UprightFont=*rm, ItalicFont=*ti, SlantedFont=*sl,
  BoldFont=*bx, BoldItalicFont=*bi, BoldSlantedFont=*bl]
\newXeTeXintercharclass\germain@greek
\def\germain@greekrange#1#2{%
  \@tempcnta=#1\relax
  \loop
    \XeTeXcharclass\@tempcnta=\germain@greek
    \ifnum\@tempcnta<#2\advance\@tempcnta\@ne\repeat}
\germain@greekrange{"0370}{"03FF}
\germain@greekrange{"1F00}{"1FFF}
\def\germain@greekprev{\rmdefault}
\def\germain@greekin{%
  \edef\germain@greekprev{\f@family}\fontfamily{germaingreek}\selectfont}
\def\germain@greekout{\fontfamily{\germain@greekprev}\selectfont}
\def\germain@greekjoin{%
  \XeTeXinterchartokenstate=\@ne
  \@tempcnta=\z@
  \loop
    \ifnum\@tempcnta=\germain@greek\else
      \XeTeXinterchartoks\germain@greek\@tempcnta=\expandafter{%
        \expandafter\germain@greekout\the\XeTeXinterchartoks\z@\@tempcnta}%
      \XeTeXinterchartoks\@tempcnta\germain@greek=\expandafter{%
        \the\XeTeXinterchartoks\@tempcnta\z@\germain@greekin}%
    \fi
    \ifnum\@tempcnta<4095\advance\@tempcnta\@ne\repeat}
% babel-french sets its own classes, and saves and restores the interchar
% state, each time a language is selected: join again after it does.
\addto\extrasfrench{\germain@greekjoin}
\addto\extrasenglish{\germain@greekjoin}
\AtBeginDocument{\germain@greekjoin}
\makeatother

\usepackage{geometry}
\geometry{a4paper,margin=25mm,marginparwidth=28mm}
\usepackage{marginnote}
\usepackage{xcolor}
\usepackage[normalem]{ulem}
\setlength{\emergencystretch}{3em}
\usepackage{hyperref}
\hypersetup{colorlinks=true,linkcolor=black,urlcolor=black,pdfborder={0 0 0}}

% --- Batch metadata ---------------------------------------------------------
% Taken as they stand by the rendering script, which shows them at the head of
% the reading view. They fix what the file claims to cover. The macro names are
% the parent project's, so that the scripts stay shared; \folder here means the
% volume's slug — `fr-9115` — and \pages counts Gallica views.

\newcommand{\folder}[1]{\gdef\grothFolder{#1}}
\newcommand{\batch}[1]{\gdef\grothBatch{#1}}
\newcommand{\pages}[2]{\gdef\grothFirst{#1}\gdef\grothLast{#2}}
\newcommand{\foldertitle}[1]{\gdef\grothFoldertitle{#1}}

% The holder's own shelfmark, printed as they write it: « Français 9115 ».
\newcommand{\shelfmark}[1]{\gdef\grothShelfmark{#1}}

% The dating, copied verbatim from the holder's catalogue — never inferred
% here. « XIXe siècle » is the BnF's; a pass that dates a leaf from its content
% says so in a \note{}, as its own inference.
\newcommand{\dating}[1]{\gdef\grothDating{#1}}

% The status notice, printed in the title block and repeated as a diagonal
% watermark on every page of the PDF. The manuscripts are in the public
% domain; what this line declares is not a right but a quality — an
% unverified first pass by a machine. The skills write the line; a file
% without it gets no watermark, so leaving it out is visible in the output.
\newcommand{\watermark}[1]{\gdef\grothWatermark{#1}}

\gdef\grothFolder{?}\gdef\grothBatch{}\gdef\grothFirst{?}\gdef\grothLast{?}%
\gdef\grothFoldertitle{}\gdef\grothShelfmark{}\gdef\grothDating{}\gdef\grothWatermark{}

% --- The résumé -------------------------------------------------------------
\newenvironment{resume}
  {\small\setlength{\parskip}{0.45em}}
  {\par\medskip\hrule\medskip}

% --- Keywords ---------------------------------------------------------------
% In English, closing the modernised reading's résumé: the modern vocabulary
% under which to search for what the leaves construct. The manifest extracts
% them to tag the volume — this line is the single source of the tags.
\newcommand{\keywords}[1]{\par\smallskip\noindent
  {\footnotesize\textsc{Keywords} --- \textit{#1}}\par}

% --- The critical apparatus -------------------------------------------------

% The Gallica view number, in the margin. This is the anchor that lets a
% disagreement over a formula be settled by looking at *one* image rather than
% a volume of 750. The site uses it to turn the facsimile as the transcription
% is scrolled. A view is one scanned image: usually one side of a leaf.
\newcommand{\page}[1]{%
  \marginnote{\footnotesize\color{gray!70}\textsf{v.\,#1}}%
  \label{page:#1}\ignorespaces}

% The BnF's pencilled foliation, where the leaf carries it and a pass can read
% it: « 198r ». This is what the literature cites — « ff. 198r–208v » — and
% the one line that turns a citation into a view. Recorded, never inferred:
% a folio a pass cannot read is not written.
\newcommand{\folio}[1]{%
  \marginnote{\footnotesize\color{gray!50}\textsf{f.\,#1}}\ignorespaces}

% The modernised reading's coarser counterpart to \page{}: which views a
% section is drawn from.
\newcommand{\pagerange}[2]{%
  \marginnote{\footnotesize\color{gray!70}\textsf{v.\,#1--#2}}%
  \ignorespaces}

% Illegible: never guessed.
\newcommand{\ill}{\textcolor{red!60!black}{[\,\ldots\,]}}

% A doubtful reading: the word is given, and flagged as uncertain.
\newcommand{\uncertain}[1]{\uline{#1}}

% An editorial addition — a missing letter, an implied word.
\newcommand{\add}[1]{\textcolor{blue!50!black}{[#1]}}

% Struck out by the writer. What was crossed out is part of the document.
\newcommand{\struck}[1]{\sout{#1}}

% The transcriber's note, on the state of the page or the mathematics.
\newcommand{\note}[1]{\par\smallskip{\small\color{gray!60!black}\textit{[#1]}}\par\smallskip}

% A marginal note of hers, distinct from the body of the text.
\newcommand{\marginal}[1]{\par\smallskip\noindent\hspace{1em}%
  {\small\color{gray!60!black}#1}\par\smallskip}

% --- Title ------------------------------------------------------------------
% The title block is French, like the documents themselves.

\AddToHook{shipout/background}{%
  \ifx\grothWatermark\empty\else
    \begin{tikzpicture}[remember picture, overlay]
      \node[rotate=52, scale=3.4, align=center, text=gray!18,
            font=\sffamily\bfseries]
        at (current page.center) {\grothWatermark};
    \end{tikzpicture}%
  \fi}

\AtBeginDocument{%
  \begin{center}
    {\large\bfseries Manuscrits de mathématiciens (Gallica) ---
      \ifx\grothShelfmark\empty\grothFolder\else\grothShelfmark\fi}\\[2pt]
    {\itshape\grothFoldertitle}\\[4pt]
    {\small vues \grothFirst--\grothLast
      \ifx\grothBatch\empty\else\ (lot \grothBatch)\fi}\\[2pt]
    \ifx\grothDating\empty\else
      {\small\color{gray!60!black}Datation du catalogue : \grothDating}\\[2pt]
    \fi
    \ifx\grothWatermark\empty\else
      {\small\bfseries\color{red!45!gray}\grothWatermark}\\[2pt]
    \fi
    {\footnotesize\color{gray!60!black}Transcription établie d'après les images IIIF de Gallica.
     Source gallica.bnf.fr / Bibliothèque nationale de France.\\
     Les lectures incertaines sont soulignées, les illisibles notées [\,\ldots\,],
     les ajouts entre crochets ; \textsf{v.} numérote les vues de Gallica,
     \textsf{f.} la foliotation au crayon de la BnF.}
  \end{center}
  \vspace{1em}}

\endinput


\folder{naf-5166}
\shelfmark{NAF 5166}
\batch{3}
\pages{41}{60}
\dating{1701-1900}
\watermark{Lecture automatique\\non vérifiée}
\foldertitle{Papiers du géomètre J.-L. LAGRANGE (1813).}

\begin{document}

\page{41}
\note{Feuillet de papier ancien, le bord supérieur déchiré, monté sur onglet. En haut à droite, à l'encre, « 12. », en chiffres appuyés suivis d'un point ; aucun folio n'est porté, voir l'en-tête. Au milieu de la page, à droite du texte, le cachet rond « BIBLIOTHÈQUE NATIONALE / R.F. / MANUSCRITS ». La page commence au milieu d'une phrase : le texte a commencé avant la vue 41.}

à l'equation que nous venons de trouver, car l'integrale de $\rho\, d\mu$ etant $-\frac{1}{ar\rho} = -\frac{\sqrt{(r^2 - 2ar\mu + a^2)}}{ar}$ il n'en peut resulter aucun terme où le facteur $r - a$, qui devient nul lorsque $r = a$, soit detruit par la division.
\note{« $\rho$ » : une lettre bouclée à longue descendante, sans la haste du \emph{p} de « peut » ou de « point » ; lue $\rho$ ici et plus bas, la lecture \emph{p} n'est pas exclue. Le radical est le grand \emph{V} suivi d'une parenthèse, comme aux vues 17 et 19. « soit » est récrit sur un premier mot.}

10. Maintenant puisque $V = \frac{4\pi a^3}{3r} + u$ (n.\textsuperscript{o} 2) on aura
\[ \frac{V}{2} + \frac{a\, dV}{dr} = \frac{2\pi a^3}{3r} - \frac{4\pi a^4}{3r^2} + \frac{u}{2} + \frac{a\, du}{dr} \]
faisant $r = a(1+y)$ pour avoir l'attraction sur un point à la surface du sphéroïde on aura
\[ \frac{V}{2} + \frac{a\, dV}{dr} = -\frac{2\pi a^2}{3} + 2\pi a^2 y + \frac{u}{2} + \frac{a\, du}{dr} \]
en negligeant les $y^2$. Mais nous venons de trouver $\frac{u}{2} + \frac{a\, du}{dr} = -2a^2\pi y$ ; ainsi on a aux quantités du second ordre près
\[ \frac{V}{2} + \frac{a\, dV}{dr} = -\frac{2\pi a^2}{3} \]
\marginal{Comme M. Laplace l'a trouvé.}
\struck{C'est l'equation de M. Laplace qui se trouve ainsi demontrée rigoureusement et mise à l'abri de \uncertain{toute} difficulté.} Au reste il est bien remarquable que le terme $2\pi a^2 y$ du à l'action de la sphere sur un point élevé au dessus de sa surface de la quantité $ay$,
\marginal{et par lequel l'equation que nous avons trouvée (n.\textsuperscript{o} 2) differoit de celle de M. Laplace,}
se trouve precisement detruit par la valeur de l'integrale de la differentielle $a^3\left(\frac{\rho}{2} + \frac{a\, d\rho}{dr}\right) y'\, d\varpi'\, d\theta' \sin\theta'$, (n.\textsuperscript{o} 4) dans laquelle le facteur $\frac{\rho}{2} + \frac{a\, d\rho}{dr}$ est toujours identiquement nul.
\note{« $\pi$ » : la constante de $\frac{4\pi a^3}{3r}$ est écrite d'un grand signe en forme de $\Pi$, dont le premier jambage est bouclé ; l'angle de l'intégrale, plus bas, est un petit $\pi$ au trait recourbé, rendu $\varpi'$ comme aux vues 17 et 19. Dans la première équation détachée, l'exposant de $a$ au premier terme est un 3 suivi d'un chiffre noyé dans une tache, et celui du second terme, 4, est récrit sur un autre chiffre. « (n.\textsuperscript{o} 2) » : le chiffre se lit 2, comme dans la note marginale. Les deux notes marginales sont appelées chacune par un signe d'insertion : la première devant la ligne biffée qui suit l'équation, la seconde après « de la quantité $ay$, ». Le passage biffé tient trois lignes, rayées chacune d'un trait continu ; « toute » est récrit et en partie couvert d'une tache. Le calcul se vérifie : $\frac{2\pi a^3}{3r} - \frac{4\pi a^4}{3r^2}$ devient $-\frac{2\pi a^2}{3} + 2\pi a^2 y$ pour $r = a(1+y)$ au premier ordre. Après « nul. », un petit signe barré ; au-dessous, un long trait horizontal ferme le texte, et le quart inférieur de la page est blanc.}

\page{43}
\note{Feuillet plus petit que les précédents, monté dans la moitié inférieure d'une feuille de garde moderne. En haut à droite, à l'encre, « 13 » : le 1 porte un long trait d'attaque qui le fait ressembler à un 7, le 3 a une queue bouclée ; il est lu 13 dans la suite 12, 13, 14…\ de ces vues. Aucun folio n'est porté, voir l'en-tête. Le feuillet ne porte qu'une ligne, au premier tiers, d'une plume fine, d'une écriture posée et penchée, aux majuscules ornées, qui n'est pas celle du texte de la vue 41 ; « M\textsuperscript{r} » est souligné. Le reste du feuillet est blanc. C'est apparemment un titre de chemise pour les feuillets qui suivent ; ce qu'il couvre au juste ne se lit pas sur le feuillet.}

Notes de M\textsuperscript{r} Dela Grange
\note{« M\textsuperscript{r} » : un M suivi d'un petit \emph{r} surélevé, souligné d'un trait ; c'est « Monsieur », non « Mademoiselle ». « Dela Grange » en deux mots, le D et le G ornés.}

\page{45}
\note{Feuillet de papier ancien, le bord supérieur déchiré, monté sur onglet. En haut à droite, à l'encre, « 14 » ; aucun folio n'est porté, voir l'en-tête. Le cachet rond de la Bibliothèque nationale est posé au milieu de la page, sur la fin de la formule de $r$. La note commence en tête de page, sans titre, sur un sujet sans lien avec le texte de la vue 41 ; elle finit sur « même axe. » vers les deux tiers de la page, et le reste du feuillet est blanc. Le $z$ est écrit à queue, comme un ʒ ; il est rendu $z$.}

\[ Ax + By + Cz = 0 \qquad (x-a)^2 + y^2 + z^2 = r^2 \]
et diff.
\[ (x-a)dx + y\,dy + z\,dz = 0 \qquad z = -\frac{Ax + By}{C} \]
donc
\[ \left(x - a - \frac{zA}{C}\right)dx + \left(y - \frac{zB}{C}\right)dy = 0 \]
\[ x - a - \frac{zA}{C} = 0 \qquad y - \frac{zB}{C} = 0 \]
\[ x = a + \frac{zA}{C}, \quad y = \frac{zB}{C} \]
donc
\[ Aa + \left(\frac{A^2}{C} + \frac{B^2}{C} + C\right)z = 0 \qquad z = -\frac{ACa}{A^2 + B^2 + C^2} \]
\[ y = -\frac{ABa}{A^2 + B^2 + C^2}, \quad x = a - \frac{A^2 a}{A^2 + B^2 + C^2} \]
donc
\[ r^2 = \frac{a^2(A^4 + A^2C^2 + A^2B^2)}{(A^2 + B^2 + C^2)^2}, \quad r = \frac{Aa}{\sqrt{A^2 + B^2 + C^2}} \]
\note{« et diff. » : « et » suivi d'une abréviation dont la dernière lettre est chargée d'un paraphe, lue « diff. ». Dans « $y^2$ » de la première ligne, le $y$ est récrit sur une parenthèse ou une autre lettre. L'exposant de $a$ dans « $A^2 a$ » est suivi d'une petite tache. Les exposants 2 sont écrits comme de petits accents circonflexes, ainsi qu'aux vues 17--19.}

Or $\frac{r}{a}$ est le sinus de l'angle que le plan fait avec l'axe des $x$ donc ce sinus est $\frac{A}{\sqrt{A^2 + B^2 + C^2}}$ \uncertain{donc} cette quantité est le cosinus de l'angle que la perpendiculaire au plan fait avec le même axe.
\note{« des $x$ » : « dy » suivi d'un $x$ isolé, lu comme l'article et la lettre. Le premier mot de la dernière phrase, un \emph{d} bouclé suivi de trois jambages, est lu avec doute « donc » ; au-dessus de lui, dans la marge, un petit signe en forme d'accent circonflexe, sans texte qui lui réponde.}

\page{47}
\note{Demi-feuillet de papier ancien, monté en haut d'une feuille de garde moderne. En haut à droite, à l'encre, « 15 », le 5 à longue queue descendante, et à sa droite, d'une plume plus appuyée, une grande capitale « A » ; aucun folio n'est porté, voir l'en-tête. Au milieu, le cachet rond de la Bibliothèque nationale, posé sur la fin de la troisième formule. En tête, au milieu, « page 623 », qui renvoie à une page que le feuillet ne nomme pas. Le calcul occupe tout le demi-feuillet et ne se rattache pas à la note de la vue 45.}

page 623

Je mets $x$ à la place de \uncertain{$b$} je differentie
\[ \left( \frac{1}{2}(1+x)^{-\frac{3}{2}} + \frac{\frac{1}{2}\frac{1}{\sqrt{1+x}}}{1 + \sqrt{1+x}} - \frac{1}{2}\frac{1}{x} \right) dx \]
\[ \frac{1}{1 + \sqrt{1+x}} = \frac{1 - \sqrt{1+x}}{-x} \]
donc on aura
\[ \frac{1}{2}dx\left( \frac{1}{(1+x)\sqrt{1+x}} - \frac{1}{x\sqrt{1+x}} + \frac{1}{x} - \frac{1}{x} \right) \]
\[ = \frac{\frac{1}{2}dx}{\sqrt{1+x}}\left( \frac{1}{1+x} - \frac{1}{x} \right) = \frac{-\frac{1}{2}dx}{x(1+x)\sqrt{1+x}} \]
\[ = -\frac{1}{2}\frac{(1+x)^{-\frac{3}{2}}}{x}dx \]
\[ = -\frac{1}{2}\frac{\left(1 - \frac{3}{2}x + \frac{3.5}{2.4}x^2 - \frac{3.5.7}{2.4.6}x^3 \text{ \&c.}\right.}{x} \]
\[ = -\frac{1}{2}\left( \frac{1}{x} - \frac{3}{2} + \frac{3.5}{2.4}x - \frac{3.5.7}{2.4.6}x^2 \right) dx \]
integrant
\[ -\frac{1}{2}lx + \frac{1}{2}\left( \frac{3x}{2} - \frac{3.5}{2.4}.\frac{x^2}{2} + \frac{3.5.7}{2.4.6}.\frac{x^3}{3} \text{ \&c.} \right. \]
+ \struck{const.} $K$, \struck{const. = \ill{}}
\note{« $b$ » : une lettre à haste et à panse fermée, qui ressemble aussi au 6 des dénominateurs ; lue $b$ avec doute. Les points des produits $3.5$, $2.4.6$ sont ceux du manuscrit. Dans la quatrième ligne de calcul, la barre de la fraction $\frac{\frac{1}{2}dx}{\sqrt{1+x}}$ est traversée d'un trait oblique, qui ne paraît pas biffer la formule. La parenthèse ouverte avant $1 - \frac{3}{2}x$ et celle de la ligne « integrant » ne sont pas fermées. Le calcul se vérifie : la somme des trois termes de la première ligne vaut bien $-\frac{1}{2}\frac{dx}{x(1+x)\sqrt{1+x}}$. À la dernière ligne, « $K$ » est écrit au-dessus de « const. » biffé ; après la virgule, « const. = » suivi d'un mot court, le tout biffé d'un trait, dont la fin ne se lit pas. Le bas du demi-feuillet est blanc.}

\page{49}
\note{Demi-feuillet de papier ancien, monté dans la moitié inférieure d'une feuille de garde moderne. En haut à droite, à l'encre, « 16 », et à sa droite une capitale « B », qui répond au « A » de la vue 47 ; aucun folio n'est porté, voir l'en-tête. Le cachet rond de la Bibliothèque nationale est posé sur la fin de la formule de la septième ligne. Le texte suit le calcul de la vue 47, dont il détermine la constante $K$ ; il remplit le demi-feuillet jusqu'au pied.}

Ayant differentié et ensuite réintegré il s'ensuit que cette integrale peut differer de la quantité proposée d'une constante qu'il s'agit de trouver.

Pour cela je fais dans la quantité donnée $b = 0$ ou $b$ infiniment petite elle devient
\[ -1 + l\frac{2}{\sqrt{b}} = -1 + l2 - \frac{1}{2}lb \]
Or dans l'integrale trouvée je fais de même $x = 0$ ou infiniment petite elle devient $-\frac{1}{2}lx + K$

Comparant et faisant $x = b$ on a $-1 + l2 = K$ donc la serie cherchée est en changeant $x$ en $b$
\[ -1 + l2 - \frac{1}{2}lb + \frac{3b}{4} - \frac{3.5}{2.4}.\frac{b^2}{4} + \frac{3.5.7}{2.4.6}.\frac{b^4}{6} \]
où $l2 - \frac{1}{2}lb = l\frac{2}{\sqrt{b}}$
\note{« réintegré » : le mot commence par un \emph{r} mal formé, lu ainsi. Dans « $-1 + l2 - \frac{1}{2}lb$ » de la première formule, le 2 est récrit sur un autre chiffre. « $l$ » est le logarithme, écrit d'un \emph{l} bouclé. Les exposants sont écrits sur la ligne, à la même hauteur que la lettre : « b2 », « b4 ». Le dernier terme porte « b4 », en chiffres nets ; le calcul de la vue 47, où le terme correspondant est en $\frac{x^3}{3}$, donnerait $b^3$. Le manuscrit est transcrit tel qu'il est. La série n'est pas close par « \&c. ».}

\page{51}
\note{Demi-feuillet de papier ancien, monté en haut d'une feuille de garde moderne. En haut à droite, à l'encre, « 17 » ; aucun folio n'est porté, voir l'en-tête. Le cachet rond de la Bibliothèque nationale est posé à droite des formules du milieu. La page commence sur une expression, sans phrase qui la rattache au feuillet « 16 » (vue 49) ; le texte s'arrête sur la formule générale, vers le bas du demi-feuillet, et le reste est blanc. L'angle est écrit d'une lettre bouclée, rendue $\varphi$ ; le module et l'amplitude sont notés $E(c, \varphi)$.}

$\sqrt{(dx^2 + dy^2)}$ soit le demi-grand axe $= a$ le demi petit axe $b$, l'excentricité sera $\sqrt{a^2 - b^2} = c$.
\[ x = a \sin\varphi \qquad y = b \cos\varphi \]
\[ dx^2 = a \cos\varphi\, d\varphi \qquad dy = -b \sin\varphi\, d\varphi \]
\[ dx^2 + dy^2 = (a^2 \cos\varphi^2 + b^2 \sin\varphi^2)\, d\varphi^2 \]
en mettant pour $\cos\varphi^2$, $1 - \sin\varphi^2$
\[ dx^2 + dy^2 = (a^2 - (a^2 - b^2)\sin\varphi^2)\, d\varphi^2 \]
\[ = (a^2 - c^2 \sin\varphi^2)\, d\varphi^2 \]
\[ = a^2\left(1 - \frac{c^2}{a^2}\sin\varphi^2\right) d\varphi^2 \]
\[ \sqrt{dx^2 + dy^2} = a\sqrt{1 - \frac{c^2}{a^2}\sin\varphi^2}\,.\,d\varphi \]
si $a = 1$ on a $\int\sqrt{1 - c^2\sin\varphi^2}\,.\,d\varphi = E(c, \varphi)$ \&

donc en general on aura
\[ \int\sqrt{dx^2 + dy^2} = aE\left(\frac{c}{a}, \varphi\right) \]
\note{« $dx^2 = a \cos\varphi\, d\varphi$ » : le 2 est écrit après $dx$, au membre gauche, alors que le membre droit est la différentielle première ; transcrit tel quel. « $\varphi$ » : une lettre faite d'une boucle ouverte suivie d'une seconde boucle, qu'on pourrait lire \emph{œ} ; la dernière formule, où elle est mieux formée, la montre comme un $\varphi$. Les carrés $\sin\varphi^2$, $\cos\varphi^2$ sont écrits ainsi, l'exposant après l'angle. Le radical de la première ligne est le grand \emph{V} suivi d'une parenthèse ; ceux des lignes suivantes ont une barre. Le « $E$ » porte au-dessus un petit trait bouclé, qui paraît un ornement de la majuscule. Après $E(c, \varphi)$, un petit signe qui ressemble à « \& ».}

\page{53}
\note{Feuillet de papier ancien, plus petit, monté dans la moitié inférieure d'une feuille de garde moderne. En haut à droite, à l'encre, « 18 » ; aucun folio n'est porté, voir l'en-tête. Le cachet rond de la Bibliothèque nationale est posé sur la seconde équation différentielle. \emph{Une autre main} que celle des vues 41--51 : une écriture soignée et arrondie, d'une plume fine, aux lettres souvent détachées, au \emph{d} à haste droite (non bouclée), au \emph{z} en forme de 2 ; les accents sont notés. La plume fine rappelle celle du titre de la vue 43, sans que le rapprochement puisse être établi sur si peu de mots. La page reprend, avec plus de phrases, le calcul de la note du feuillet « 14 » (vue 45) : mêmes équations, même dernière phrase à peu de mots près. Le texte finit vers les deux tiers du feuillet ; le reste est blanc.}

\[ z = -\frac{Ax + By}{C} \]
\[ dz = -\frac{A\,dx + B\,dy}{C} \]
Substituant cette valeur de $dz$ dans l'expression differentielle de la distance
\[ (x-a)dx + y\,dy + z\,dz = 0 \]
elle devient
\[ \left(x - a - \frac{zA}{C}\right)dx + \left(y - \frac{zB}{C}\right)dy = 0 \]
le resultat de ce calcul \struck{est de} donne pour l'equation du plan passant par l'origine des coordonnées et perpendiculaire à la droite pour laquelle les coordonnées seroient proportionnelles aux quantités $A$, $B$, $C$
\[ \frac{A}{\sqrt{A^2 + B^2 + C^2}}x + \frac{B}{\sqrt{A^2 + B^2 + C^2}}y + \frac{C}{\sqrt{A^2 + B^2 + C^2}}z = 0 \]
$\frac{r}{a}$ est le sinus de l'angle que le plan fait avec l'axe des $x$.
\note{« proportionnelles » est coupé « proporti-/onnelles » en fin de ligne, et sa finale est récrite. « est de » est biffé d'un trait, « donne » écrit à la suite. Le $z$ de ces formules est écrit comme un 2 à queue ; le contexte, et la vue 45, le font lire $z$. La barre du dernier radical est doublée d'un second trait.}

\page{55}
\note{Une même feuille de garde moderne porte deux pièces de papier ancien, collées l'une au-dessus de l'autre ; chacune a son numéro à l'encre en haut à droite, « 19 » sur la pièce du haut, « 20 » sur celle du bas ; aucun folio n'est porté, voir l'en-tête. Un cachet rond de la Bibliothèque nationale est posé sur chacune.}

\note{Pièce du haut, « 19 » : un demi-feuillet écrit en largeur, sur sa moitié supérieure, de la main soignée et arrondie de la vue 53, aux chiffres moulés : le 2 en forme de \emph{z}, le 5 en forme de \emph{ſ}. En tête à gauche, « Pag. 623 » souligné, le même renvoi que « page 623 » en tête de la vue 47 ; la pièce reprend en effet la série en $b$ des vues 47 et 49, par un autre chemin. Le reste du demi-feuillet est blanc.}

Pag. 623

A cause de $\log.\frac{1 + \sqrt{(1+b)}}{\sqrt{b}} = \log.(1 + \sqrt{(1+b)}) - \log.\sqrt{b}$ il faut que l'on ait
\[ -\frac{1}{\sqrt{(1+b)}} + \log.(1 + \sqrt{(1+b)}) = \log.2 - 1 + \frac{3}{2}.\frac{b}{2} - \frac{3.5}{2.4}.\frac{b^2}{4} + \text{\&c.} \]
\[ \log.(1 + \sqrt{(1+b)}) = \sqrt{(1+b)} - \frac{1+b}{2} + \frac{(1+b)^{\frac{3}{2}}}{3} - \frac{(1+b)^2}{4} + \frac{(1+b)^{\frac{5}{2}}}{5} - \text{\&c.} \]
\[ = 1 + \frac{b}{2} - \frac{b^2}{8} + \text{\&c.} - \frac{1+b}{2} + \frac{1 + \frac{3}{2}b + \frac{3}{8}b^2 - \text{\&c.}}{3} \]
\[ - \frac{1 + 2b + b^2}{4} + \frac{1 + \frac{5}{2}b + \frac{15}{8}b^2 + \text{\&c.}}{5} - \text{\&c.} \]
\[ -\frac{1}{\sqrt{(1+b)}} = -(1+b)^{-\frac{1}{2}} = -1 + \frac{1}{2}b - \frac{1.3}{4.2}b^2 + \frac{1.3.5}{2.4.6}b^3 - \text{\&c.} \]
\[ \log.2 = 1 - \frac{1}{2} + \frac{1}{3} - \frac{1}{4} + \frac{1}{5} - \text{\&c.} \]
\note{Le « A » initial est une capitale ornée qu'on prendrait pour un C. Le dénominateur 3 de la troisième ligne est récrit en gras sur un autre chiffre. Le signe $-$ devant $\frac{1 + 2b + b^2}{4}$ est une tache allongée, repassée. Dans « $\frac{1.3}{4.2}$ », le dénominateur, en partie sous le cachet, se lit « 4.2 » et non « 2.4 » comme aux autres termes ; lecture douteuse. Les points des produits sont ceux du manuscrit ; « log. » est abrégé avec un point.}

\note{Pièce du bas, « 20 » : un morceau de papier déchiré, couvert de calculs numériques à l'encre, sans un mot de texte, d'une écriture rapide où le 4 est fait comme un A (« 2A », « 13 8A »). On y lit notamment, en haut à gauche, la soustraction « 7 25 / 1 39 / 5 46 » ; la multiplication « 3 4 6 » par « 4 $\frac{1}{2}$ », avec les produits « 13 8 4 », « 1 7 3 » et la somme « 1 5 5 7 » ; au milieu, une colonne « 60$'$ 0$''$ / 10$'$ 14 / 16 0 », et à droite l'addition « 7681 / 5740 / 13421 » ; plus bas, deux petits calculs en degrés et minutes, « 30 \ill{}4 / 2 44 / 27 20 / 51 / 28 11 », une tache d'encre couvrant le chiffre avant le 4, le 28 biffé et « 26 29 » écrit dessous ; à droite de chacun des deux calculs, un double croissant barré d'un trait, qui a la forme du signe de la lune ; une ligne « 5 14 29 57 » et, à sa droite, « 26 | 10 4 ||| 17 0 9 | 26 » ; puis un tableau de logarithmes :}
\begin{quote}
SC. 5 6 56 — 1999658\ill{}0

SC. 7 10 48 — 99982668

ST.

S.DAB 84 57 18 — 99983142

84 55 35

1 43
\end{quote}
\note{Dans ce tableau, deux traits obliques croisés relient « SC. 5 6 56 » et « SC. 7 10 48 » aux deux premiers logarithmes, en sens inverse ; « ST. » est suivi d'un trait long sans nombre. Les tirets de cette transcription tiennent la place de l'espace qui sépare, sur la pièce, l'angle de son logarithme ; les nombres commençant par 9999 ont la forme de logarithmes tabulaires de sinus ; « SC. » et « S.DAB » paraissent abréger « sinus complément » et « sinus de l'angle DAB », ce qui est une lecture de ce passage, non du feuillet. Au premier nombre, un petit chiffre entre le 8 et le 0 final, sous la croix des traits, ne se lit pas (1 ou 2). Sous les deux derniers chiffres du troisième, « 42 », est écrit « 46 », souligné. La soustraction « 84 57 18 − 84 55 35 = 1 43 » est juste. Le reste de la pièce est un enchevêtrement de retenues et de reports qui n'est pas transcrit. En bas à droite, retourné tête-bêche, un petit groupe « 100 | 52 / 1001 », lecture douteuse. Des chiffres et des lettres au crayon (« G », « I », « D », « 30 », « 40 », « 150 ») se voient à l'envers entre les colonnes : c'est le crayon du verso (vue 56) vu par transparence.}

\page{56}
\note{Revers des deux pièces de la vue 55. Celle du haut (dos de la pièce « 19 ») est blanche ; on n'y voit que le texte du recto par transparence. Celle du bas (dos de la pièce « 20 ») est écrite : des tableaux de chiffres au crayon, pâles, dans des cases tracées au crayon, les uns à l'endroit, les autres tête-bêche, traversés de longues lignes obliques, et un petit groupe à l'encre écrit en travers. Aucun numéro. Le 4 y est fait comme un A, comme au recto. Seul ce qui se lit sûrement est donné ici ; le reste, à demi effacé et croisé de traits, n'est pas transcrit. À l'endroit, en haut : une colonne « 54 / 20 4 / 46 / 11 4 / 48.42 », puis « 27 1 37 » et « 437.38 » ; dessous, une case à trois lignes numérotées « 50 / 51 / 52 », avec en regard « 0 43 33 », « 47 30 3\ill{} », « 29 28 ». À l'encre, en travers, à droite de ces cases : « 140 | 60 / 7 || 7. », et au-dessous « 420 », souligné. Tête-bêche, dans la moitié inférieure : un tableau dont la première colonne croît par demi-unités, « 55.30 / 56 0 / 56.30 / 57.0 », puis « 57.30 », « 58 0 » et deux lignes moins sûres ; en regard, des colonnes de minutes et de fractions en demis, et, dans les cases, des capitales isolées, dont « G », « K », « E » et « D ». Près du bord, un croissant au crayon, le même signe qu'au recto.}

\page{58}
\note{Grand feuillet de papier ancien, monté dans la moitié inférieure d'une feuille de garde moderne. En haut à droite, à l'encre, « 21 » ; aucun folio n'est porté, voir l'en-tête. Au milieu, le cachet rond de la Bibliothèque nationale, posé sur « les doubles integrales ». \emph{Une autre main} que celle des vues 41--51 : une cursive régulière, penchée à droite, à pleins et déliés marqués, aux mots souvent prolongés par un long trait horizontal ; le \emph{d} des différentielles a une haste droite ; beaucoup de corrections dans l'interligne. Elle ressemble de près à la main soignée des vues 53 et 55 (mêmes traits de sortie horizontaux, même \emph{d} droit, même 5 en forme de \emph{ſ}), écrite ici plus vite ; que ce soit la même main n'est pas établi. C'est un texte à la première personne, adressé à un « vous », sans formule d'appel ni date : une lettre ou le brouillon d'une lettre. Il commence en tête du feuillet et se poursuit au verso (vue 59) au milieu d'une phrase. Le papier laisse voir par transparence, à l'envers, l'écriture du verso.}

Je ne suis pas fort étonné du resultat que vous m'annoncez, car je n'avois guere de confiance dans mon travail j'y avois été entraîné par une analogie qui me sembloit frappante mais que je ne me croyois pas trop en état d'apprecier, je vous ai la plus grande obligation des soins que vous avez pris pour me faire \struck{connoître} obtenir un jugement qui m'éclaire \struck{entierement} sur les fautes que j'ai commises
\note{« étonné », « entraîné » : les deux participes finissent par un \emph{é} suivi d'un court trait de sortie ; ce trait pourrait porter un second \emph{e}, sans qu'on puisse l'en distinguer. Le genre de celui qui écrit ne se lit donc pas sur ces deux mots ; voir plus bas « egarée » et « persuadée », où la finale \emph{-ée} est nette. « connoître » et « entierement » sont biffés chacun d'un trait ; lecture de « connoître » sous le trait incertaine.}

Ce qui me surprend le plus, c'est l'equation que vous m'indiquez d'après M\textsuperscript{r} De la Grange ; si je ne savois d'où elle vient et par conséquent la confiance qu'elle merite, je l'aurois crue au premier coup d'œil formée de la somme des deux equations
\[ p\frac{d^2z}{dt^2} + 2k^2\left(\frac{d^4z}{dx^4} + \frac{d^4z}{dx^2dy^2}\right) = 0 \]
\[ p\frac{d^2z}{dt^2} + 2k^2\left(\frac{d^4z}{dy^4} + \frac{d^4z}{dy^2dx^2}\right) = 0, \]
mais je sens bien qu'il faut que je me trompe. Si je me suis egarée dans les doubles integrales je ne dois pas m'étonner que l'équation qu'a trouvée M\textsuperscript{r} De la Grange ne monte qu'au quatrième degré je suis très persuadée de la realité de mon erreur puisqu'elle a été reconnue par des juges dont je respecte l'opinion mais je \struck{suis} ne sais en quoi elle consiste. Si l'équation
\[ \int dz\,dy \int X\,dm + \int dx\,dz\; Y\,dm - 2\int dx\,dy \int Z\,dm = 2\,\mathrm{I}\,d^2z\,.\,dx\,dy \]
est réellement déduite, comme j'ai cherché à l'établir dans la note, des formules de la mécanique analytique \uncertain{p. 53} \uncertain{1\textsuperscript{re}} édition
\note{« $p$ » en tête des deux équations : un \emph{p} traversé d'un trait oblique, ou récrit sur une autre lettre ; lecture douteuse. Dans la seconde, le 2 de « $2k^2$ » est récrit sur le signe $+$ ; au-dessus de « $=0$, », dans l'interligne, un petit « $k^2$ » isolé, sans signe d'insertion. « crue », « trouvée » : ainsi. Dans la dernière équation, le second terme n'a qu'un signe d'intégration, écrit devant $dx\,dz$ ; le $Y$ est repassé. Le signe entre « $2$ » et « $d^2z$ » est un trait vertical épais, rendu $\mathrm{I}$ faute de mieux. Le renvoi final se lit « p », un point, un chiffre fait comme un \emph{ſ} (le 5 de cette plume, voir la vue 55), un 3, puis un chiffre suivi de lettres surélevées, lu « 1\textsuperscript{re} » ; la lecture de l'ensemble est douteuse. La phrase continue au verso, vue 59.}

\page{59}
\note{La vue montre le document ouvert : à gauche, le verso du feuillet « 21 » (vue 58), écrit sur toute sa hauteur, sans numéro ; à droite, la partie gauche d'une autre page, qui est le recto du feuillet « 22 » photographié en entier à la vue 60 et transcrit là. Au-dessus, la feuille de garde. La page de gauche continue sans rupture la phrase laissée au pied de la vue 58 ; elle est de la même main, avec des corrections dans l'interligne.}

ne convient-il donc pas de la differencier quatre fois \struck{pour} savoir deux fois par rapport à $x$ et deux fois par rapport à $y$ n'est-il pas au moins incontestable qu'en fesant abstraction du second membre qui contient le terme du à l'elasticité et prenant $X = Y$ on \struck{a pour le mo} par ce procédé des differentiations successives obtient \struck{pour le mouvement} l'equation \struck{connue de la surface tendue}
\[ X\left(\frac{d^2z}{dx^2} + \frac{d^2z}{dy^2}\right) = 2Z \]
qui est au fond la même que celle donnée par M\textsuperscript{r} La Grange lui même p. \uncertain{151} de la 1\textsuperscript{ère} édition, et qui appartient à l'equilibre de la surface tendue par des forces egales suivant les \struck{directions} $x$ et les \uncertain{$z$}. comment se peut-il donc faire que ce même procédé devienne vicieux lorsqu'il s'applique au terme dependant de l'elasticité.
\note{« par ce procédé des differentiations successives » est écrit dans l'interligne, au-dessus de « obtient » et de ce qui suit, sans signe d'insertion ; il est placé ici devant « obtient ». Les deux passages biffés qui l'encadrent le sont d'un trait continu ; « a pour le mo » commence à la fin d'une ligne et finit au début de la suivante. « p. 151 » : le premier chiffre est un 1, le deuxième le 5 en forme de \emph{ſ} de cette plume, le troisième est chargé d'encre ; lecture douteuse. « 1\textsuperscript{ère} » : un 1 suivi de lettres surélevées. « $x$ » est souligné. Après « les », une lettre anguleuse à queue, qui a la forme d'un $z$ ; le sens voudrait $y$.}

Vous me faites remarquer que l'equation que M\textsuperscript{r} De la Grande a trouvée en suivant mon hypothese contient celle de la lame elastique tandis que celle que j'ai deduite de la même supposition ne jouit pas de cet avantage. je \struck{\ill{}} m'ettois déja proposé \struck{cette} la même objection mais j'avois cru y repondre en observant, que \struck{le resultat de la differenti} \struck{qu'en differentiant par rapport} à une quantité qui est constante pour la lame on
\note{« De la Grande » : ainsi, avec un \emph{d}, alors que la vue 58 et la ligne plus haut portent « Grange ». Après « je », un mot de quatre ou cinq lettres biffé d'un trait épais, illisible. « m'ettois » : ainsi, deux \emph{t} sous une seule barre. « la même » est écrit au-dessus de « cette » biffé. Au-dessus de « qu'en differentiant par rapport », biffé, l'interligne porte « que le resultat de la differenti », dont la fin est biffée elle aussi ; le « que » initial paraît échapper au trait. La phrase reste ainsi sans construction sûre. Elle continue au pied de la page sur « on » et se poursuit, selon toute apparence, en tête du feuillet « 22 » (vue 60).}

\page{60}
\note{Recto du feuillet « 22 », photographié en entier ; à gauche de la vue, le bord de la page de la vue 59. En haut à droite, à l'encre, « 22 » ; aucun folio n'est porté, voir l'en-tête. Même main qu'aux vues 58--59. La page continue la phrase laissée au pied de la vue 59. Son premier tiers est un brouillon très repris : des lignes entières biffées d'un trait, et d'autres écrites entre elles, en retrait, sans signe d'insertion. La transcription suit l'ordre des lignes sur la page, en plaçant chaque ajout de l'interligne devant les mots qu'il surmonte ; la construction de plusieurs phrases reste incertaine. Le texte finit sur « géomètre. », suivi d'un point, au deuxième tiers de la page ; le reste est blanc, et il n'y a ni formule finale ni signature.}

satisfesoit à l'équation par la supposition \uncertain{A.} de constante $= 0$ de sorte que le \struck{coëfficient} facteur qui compose l'équation propre à la \struck{tant que par conséquent il fut necessaire, pour ce \ill{} cas que} surface, n'étoit \struck{n'est seroit pas} pas necessairement egal à zéro lorsqu'il\ldots\ le \struck{facteur que est zero pour la surface} l'autre facteur \uncertain{sujet de} la lame ; que \struck{est fait zero} \struck{et par conséquent que} l'équation qui en \struck{resulteroit fut applicable} au reste j'avois cru que la formule considerée avant les differentiations successives renfermoit le \struck{étoit également applicable au} cas de la lame, et c'est ce que j'ai eu l'intention de prouver à la fin de la note.
\note{« \uncertain{A.} » : une capitale suivie d'un point, qui pourrait être un \emph{d} bouclé ; lecture douteuse. « de sorte que le coëfficient facteur qui compose l'équation propre à la » est écrit en retrait, entre la première ligne et la ligne biffée « tant que par conséquent\ldots », et se continue par « surface » au début de la ligne suivante. Entre « pour ce » et « cas », un mot court est noyé sous une rature. « n'étoit » est écrit au-dessus de « n'est seroit pas », biffés ; les deux négations se suivent ainsi. Après « lorsqu'il », la ligne s'arrête ; la suite, « le facteur\ldots », est à la ligne au-dessous. « \uncertain{sujet de} la lame ; » est écrit dans l'interligne, en retrait, sous « l'autre facteur ». « renfermoit le » est écrit au-dessus de « étoit également » et « applicable au », biffés.}

Puisque me voila revenue à cette équation trouvée par La Grange, je remarquerai qu'il est singulier que l'on puisse obtenir les deux dont j'ai deja parlé \struck{est la} et dont celle-ci semble \struck{la somme} en substituant tout bonnement dans l'équation de la lame $\frac{1}{r} + \frac{1}{r'}$ à la place de $\frac{1}{r}$ et \struck{mettant} \struck{pour l'une} \struck{\ill{}} \struck{\ill{}} $Z, x$ pour l'autre $Z, y$ à la place de $y$ et $Z$, je sais bien que ce n'est pas ainsi qu'a du operer M\textsuperscript{r} De la Grange mais cette idée m'a tellement frappée à la vue de son équation que je n'ai pû resister à la tentation de vous la communiquer en protestant au reste de ma parfaite soumission aux oracles \struck{des maîtres et} jugement d'un autre grand géomètre.
\note{« revenue », « frappée » : ainsi, la finale \emph{-ée} nette dans « frappée ». « par La Grange » : « La Grange » est au début de la ligne suivante, sans « M\textsuperscript{r} ». « et dont celle-ci semble » et « la somme », biffé, sont écrits dans l'interligne au-dessus de « est la », biffé, et de « en substituant ». « tout bonnement » est écrit au-dessus de « dans ». « pour l'une » est écrit au-dessus de « mettant », et biffé comme lui ; suivent deux mots courts biffés d'un trait épais, illisibles, dont le premier paraît fait de deux lettres de formule. Le $Z$ est la grande lettre du « $2Z$ » de la vue 59, non le petit $z$ des différentielles ; celui de « $y$ et $Z$ » est récrit. « jugement d'un autre grand » est écrit à la suite de « des maîtres et », biffés, d'une encre plus pâle et d'un trait plus fin, montant vers la droite ; « géomètre. » est à la ligne.}

\end{document}
